% Source: https://artofproblemsolving.com/community/c3h1335086_what_are_mods
% Stable identity: aops-topic-1335086-q01
\textbf{(a)} Write $a=b+mk$ and $c=d+ml$ with integers $k,l$. Then
\[
(a+c)-(b+d)=m(k+l),\qquad ac-bd=(b+mk)(d+ml)-bd=m(bl+dk+mkl),
\]
both multiples of $m$. So $a+c\equiv b+d$ and $ac\equiv bd\pmod m$. Applying the product rule repeatedly, $a\equiv b$ implies $a^n\equiv b^n\pmod m$ for every positive integer $n$.

\textbf{(b)} $3^3=27\equiv-1\pmod7$, so $3^6=(3^3)^2\equiv1\pmod7$. Since $100=6\cdot16+4$, we get $3^{100}=(3^6)^{16}\cdot3^4\equiv1^{16}\cdot81\pmod7$, and $81=11\cdot7+4$. The remainder is $\boxed{4}$.
